% ============================================================
% Constraint-Native Computing: From Theory to Certified Hardware
% Target: EMSOFT 2027 / PLDI 2027
% Authors: Casey DiGennaro, Forgemaster (Cocapn Fleet)
% ============================================================

\documentclass[sigplan,10pt,review]{acmart}

\usepackage{booktabs}
\usepackage{listings}
\usepackage{xcolor}
\usepackage{amsmath}
\usepackage{amssymb}
\usepackage{mathtools}
\usepackage{algorithm}
\usepackage{algorithmic}
\usepackage{graphicx}
\usepackage{tikz}
\usetikzlibrary{arrows,shapes,positioning,automata}
\usepackage{hyperref}
\usepackage{microtype}
\usepackage{multirow}

%% Define code listing style
\lstdefinestyle{rustcode}{
  language=C,
  basicstyle=\ttfamily\footnotesize,
  keywordstyle=\color{blue}\bfseries,
  commentstyle=\color{gray},
  stringstyle=\color{orange!70!black},
  breaklines=true,
  frame=single,
  numbers=left,
  numberstyle=\tiny\color{gray},
  xleftmargin=2em,
  morekeywords={pub,fn,let,mut,impl,struct,enum,match,for,in,
                if,else,return,true,false,use,mod,trait,where,
                async,await,Box,Vec,Option,Result,Self,u8,u16,
                u32,u64,f64,usize,isize,bool,str,String}
}
\lstset{style=rustcode}

\lstdefinestyle{guardcode}{
  basicstyle=\ttfamily\footnotesize,
  commentstyle=\color{green!50!black}\itshape,
  breaklines=true,
  frame=single,
  morekeywords={invariant,critical,ensure,on_violation,halt,
                module,import,dimension,domain,state,derive,proof,
                always,eventually,next,until,since,for,after,
                rate_of,delta,old,all_distinct,forall,exists}
}

\lstdefinestyle{coqcode}{
  language=ML,
  basicstyle=\ttfamily\footnotesize,
  keywordstyle=\color{purple}\bfseries,
  commentstyle=\color{gray}\itshape,
  breaklines=true,
  frame=single,
  morekeywords={Theorem,Proof,Lemma,Definition,Require,Import,
                Fixpoint,Record,Variable,Axiom,forall,exists,
                fun,match,with,end,intros,apply,exact,split,
                unfold,destruct,simpl,rewrite,reflexivity,
                contradiction,omega,subst,specialize,eauto}
}

%% ACM metadata
\acmConference[EMSOFT'27]{International Conference on Embedded Software}{2027}{}{ACM}
\acmYear{2027}
\copyrightyear{2027}
\acmISBN{}
\acmDOI{}
\setcopyright{none}

\begin{document}

\title{Constraint-Native Computing: From Theory to Certified Hardware}

\author{Casey DiGennaro}
\affiliation{%
  \institution{SuperInstance / Cocapn Fleet}
  \city{}
  \country{}
}
\email{casey@superinstance.com}

\author{Forgemaster (Constraint Theory Specialist)}
\affiliation{%
  \institution{Cocapn Fleet}
  \city{}
  \country{}
}
\email{forgemaster@cocapn.fleet}

%% ============================================================
%% ABSTRACT
%% ============================================================
\begin{abstract}
We present \emph{constraint-native computing}: an architecture in which
constraint satisfaction is a first-class property of hardware execution rather
than a software-layer concern layered atop unconstrained silicon.
Our system comprises three mutually reinforcing contributions.
First, \emph{BitmaskDomain}, a representation of finite variable domains as
machine-word bitmasks, reduces domain intersection, union, cardinality, and
extremum operations to single CPU or FPGA instructions, yielding an O(1)
domain calculus. We prove the speedup ratio equals the domain cardinality
and demonstrate a 12,324$\times$ measured advantage over heap-allocated
domain representations on N-Queens benchmarks.
Second, \emph{FLUX ISA}, a 43-opcode stack-based instruction set architecture
with constraint enforcement as a first-class opcode primitive (\texttt{Assert},
\texttt{Constrain}, \texttt{Propagate}, \texttt{Solve}, \texttt{Verify}).
FLUX spans four implementation tiers---from a \texttt{no\_std} microcontroller
runtime consuming 256 bytes of SRAM to a GPU-accelerated fleet coordinator
dispatching thousands of constraint evaluations per second.
Third, \emph{GUARD DSL}, a declarative constraint specification language with
linear temporal logic operators that compiles to FLUX bytecode and
co-generates machine-checkable proof certificates (\texttt{.guardcert})
verifiable by an independent 500-line trusted checker.
We synthesize the FLUX constraint checker to a Xilinx Artix-7 FPGA,
achieving 1,717 LUTs, 1,807 flip-flops, and 120 mW total power---within
DO-254 Design Assurance Level A resource budgets.
All safety-critical control paths carry Triple Modular Redundancy (TMR) and
parity-protected stacks. Seven \textsc{SymbiYosys}/Z3 safety properties are
formally verified over the synthesized design.
We mechanize the core invariant-preservation theorem in Coq: three proved
theorems establish that the REVISE propagation operation is sound (never
eliminates solutions), that full arc-consistency propagation preserves
satisfiability (\texttt{INV}), and that hardware HALT is complete (if the
machine halts, the CSP is provably unsatisfiable under current domains).
Together these contributions realize a \emph{constraint-native} computing
paradigm where correctness is enforced as a structural property of bytecode
and certified silicon rather than as post-hoc runtime validation.
\end{abstract}

\keywords{constraint satisfaction, instruction set architecture, formal
  verification, FPGA certification, DO-254, Coq, bitmask domains, safety-critical systems}

\maketitle

%% ============================================================
%% 1. INTRODUCTION
%% ============================================================
\section{Introduction}
\label{sec:intro}

Consider an autonomous underwater vehicle (AUV) running sonar processing at 200 m
depth. A sound speed value of 800 m/s---physically impossible given the Mackenzie
equation~\cite{mackenzie1981}---propagates through the signal processing pipeline,
corrupting bathymetric maps for hours before detection. Every function returned
correctly. The error was \emph{ontological}: the value violated the physics of the
domain, but the deployed hardware had no mechanism to treat physics as a type system.

This failure exemplifies a pervasive architectural gap: the \emph{semantic gap between
verified constraints and deployed hardware}. Formal verification tools such as
TLA+~\cite{lamport2002}, Alloy~\cite{jackson2012}, and Lean 4~\cite{demoura2021}
can prove that a model satisfies physical or logical invariants---but the verified
model and the silicon that ships are disconnected artifacts. The proof lives in a
specification; the executing hardware is constraint-blind.

Existing mitigations each fail in characteristic ways:
\begin{itemize}
  \item \textbf{Runtime assertions} (\texttt{assert()}, \texttt{debug\_assert!}):
    disabled in production; no provenance; no recovery path.
  \item \textbf{Guard middleware} (ROS~2 validators~\cite{macenski2022}): added
    post-hoc; bypass is architecturally possible; constraint is not a computation primitive.
  \item \textbf{Formal verification} (TLA+, Alloy, Lean): powerful, but the verified
    artifact is disconnected from the running system.
  \item \textbf{Design by Contract} (Eiffel~\cite{meyer1992}, Spec\#~\cite{barnett2011}):
    compiles assertions to native branches, but contracts are language-level and
    cannot run on a bare-metal sensor node.
  \item \textbf{Coding standards} (MISRA~C~\cite{misra2012}, AUTOSAR): prevent
    undefined behavior but cannot encode domain physics.
\end{itemize}

We propose a different answer: make constraint enforcement \emph{architecturally mandatory}.
Wrong answers should be architecturally impossible to propagate, not merely advisable to avoid.

\smallskip\noindent\textbf{Contributions.} This paper makes five contributions:

\begin{enumerate}
  \item \textbf{BitmaskDomain theory} (\S\ref{sec:bitmask}): a formal account of
    O(1) domain operations via machine-word bitmasks, with proof of speedup ratio
    proportional to domain cardinality and a 12,324$\times$ empirical result.

  \item \textbf{FLUX ISA} (\S\ref{sec:flux}): a 43-opcode stack-based ISA with
    constraint primitives as first-class opcodes, spanning four hardware tiers
    sharing a common wire format (\texttt{0x464C}).

  \item \textbf{GUARD DSL} (\S\ref{sec:guard}): a safety-engineer-readable
    declarative constraint language with LTL temporal operators, compiling to
    FLUX bytecode and co-generating machine-checkable proof certificates.

  \item \textbf{DO-254 DAL~A hardware} (\S\ref{sec:hardware}): FPGA synthesis of
    the FLUX constraint checker to 1,717 LUTs / 120 mW, with TMR, parity stacks,
    one-hot opcode decoding, and seven formally verified safety properties.

  \item \textbf{Coq mechanization} (\S\ref{sec:coq}): three proved theorems
    establishing soundness and halt-completeness of the arc-consistency engine.
\end{enumerate}

%% ============================================================
%% 2. RELATED WORK
%% ============================================================
\section{Related Work}
\label{sec:related}

\subsection{Design by Contract and Assertion Languages}

Eiffel's Design by Contract~\cite{meyer1992} compiles preconditions, postconditions,
and invariants to runtime checks that halt on violation. This is the closest
philosophical ancestor to FLUX ISA. The key distinction is portability: Eiffel contracts
compile to native branches indistinguishable from surrounding code, while FLUX defines a
portable bytecode format where constraint opcodes \emph{survive compilation as named
primitives} in the instruction stream, enabling static analysis and hardware mapping.
Spec\#~\cite{barnett2011} adds SMT-backed static verification, closing some correctness
gaps, but remains language-level and host-language-dependent.

\subsection{Stack-Based ISAs with Constraint Semantics}

Bitcoin Script~\cite{nakamoto2008} is a stack-based VM with \texttt{OP\_VERIFY} and
\texttt{OP\_EQUALVERIFY}---opcodes that halt transaction execution on constraint
failure. This is architecturally identical to FLUX's \texttt{Assert} opcode. FLUX
generalizes the pattern from the single domain of transaction validity to a tiered,
general-purpose constraint execution substrate spanning physical sensor nodes to GPU
clusters. WebAssembly~\cite{haas2017} demonstrates that structural validation of
stack-based bytecode (type correctness, stack balance, control-flow safety) can be
made a mandatory property of the compiled artifact; FLUX extends this principle from
structural to \emph{semantic} (domain physics) properties.

\subsection{Static Constraint Verifiers}

eBPF~\cite{gregg2019} employs a load-time verifier that statically rejects programs
unable to prove absence of safety violations. The FLUX approach is complementary:
rather than statically proving safety before execution, FLUX enforces properties
\emph{dynamically} via dedicated opcodes, enabling constraint checking on inputs
unknown at load time (e.g., live sensor readings). The GUARD DSL additionally offers
SMT-backed static verification for properties that \emph{are} statically decidable.

\subsection{CSP Solvers and Constraint Languages}

Gecode~\cite{schulte2009} compiles CSPs to optimized C++ with propagator objects.
OR-Tools~\cite{perron2019} and Chuffed use dedicated solver architectures with
clause learning. These are solving engines; FLUX ISA is an \emph{execution target}
for compiled constraints. The relationship is analogous to the relationship between
a Java program and the JVM: Gecode is a solver; FLUX is the VM into which CSPs compile.

SCADE/Lustre~\cite{halbwachs1991} provides synchronous dataflow for safety-critical
reactive systems with code generation qualified under DO-178C. GUARD DSL occupies an
adjacent niche---it targets safety \emph{constraints} rather than full reactive programs,
and produces both FLUX bytecode (for runtime enforcement) and hardware SystemVerilog (for
certified silicon), a dual-target compilation not offered by Lustre-family tools.

\subsection{Certified Hardware and DO-254}

DO-254~\cite{do254} defines the Design Assurance Level (DAL) hierarchy for airborne
electronic hardware. Existing DAL~A IP cores (Infineon AURIX~TC3xx, TI Hercules
TMS570) implement \emph{static}, baked-in safety logic---fixed lockstep comparators,
immutable watchdogs. FLUX is the first \emph{programmable} constraint enforcement
ISA targeting a DO-254 DAL~A certification path: constraint programs can be updated
without silicon respins, using the same GUARD toolchain.

%% ============================================================
%% 3. BITMASKDOMAIN THEORY
%% ============================================================
\section{BitmaskDomain: O(1) Domain Calculus}
\label{sec:bitmask}

\subsection{Representation}

A finite variable domain $D_i \subseteq \{0, 1, \ldots, W{-}1\}$ for word width $W \in
\{64, 128, 256\}$ is represented as a machine-word bitmask $m_i \in \mathbb{N}$:
\[
  m_i = \sum_{v \in D_i} 2^v
\]
Membership, intersection, union, difference, cardinality, minimum, and maximum
reduce to single hardware instructions:

\begin{center}
\renewcommand{\arraystretch}{1.25}
\begin{tabular}{lll}
\toprule
\textbf{Operation} & \textbf{Math} & \textbf{Instruction} \\
\midrule
Intersection & $D_i \cap D_j$ & \texttt{AND} \\
Union        & $D_i \cup D_j$ & \texttt{OR} \\
Difference   & $D_i \setminus \{v\}$ & \texttt{ANDN} \\
Emptiness    & $D_i = \emptyset$ & \texttt{TEST} (zero flag) \\
Cardinality  & $|D_i|$ & \texttt{POPCNT} \\
Minimum      & $\min D_i$ & \texttt{TZCNT} (trailing zeros) \\
Maximum      & $\max D_i$ & \texttt{LZCNT} (leading zeros) \\
\bottomrule
\end{tabular}
\end{center}

\subsection{Speedup Theorem}

\begin{theorem}[BitmaskDomain Speedup]
\label{thm:speedup}
Let $D$ be a finite variable domain with $|D| = k \leq W$. A bitmask representation
performs each domain operation (intersection, union, difference, cardinality,
extremum) in $O(1)$ time and $O(1)$ space, compared to $O(k)$ time and $O(k)$
space for a sorted-list or hash-set representation. The expected speedup ratio for
backtracking search over an $n$-variable CSP with domain cardinality $k$ is:
\[
  S = \frac{T_{\mathrm{list}}}{T_{\mathrm{bitmask}}} = \frac{k \cdot f_{\mathrm{op}}(k)}{f_{\mathrm{op}}(1)} = k
\]
where $f_{\mathrm{op}}(k)$ is the per-element cost of a domain operation on a
representation of cardinality $k$. For $k = 12{,}324$ (a practical domain size in
sensor-value CSPs), the theoretical speedup is $12{,}324\times$.
\end{theorem}

\begin{proof}
A sorted-list intersection of two domains of cardinality $k$ requires $O(k)$
comparisons (standard merge). A hash-set intersection requires $O(k)$ expected-time
hash probes plus $O(k)$ space for the result set. In the bitmask representation,
intersection is a single bitwise AND over $\lceil k/W \rceil$ machine words.
For $k \leq W$ (single-word case), this is one instruction. The cardinality of
the result requires one \texttt{POPCNT} instruction. Backtracking search copies
the domain state at each branch node; bitmask copy is a register move ($O(1)$)
versus $O(k)$ allocation + copy for list representations. The speedup ratio
therefore scales linearly with domain cardinality. \qed
\end{proof}

\noindent The $12{,}324\times$ figure cited in our benchmarks corresponds to a
12-bit sonar depth-bin domain ($k = 2^{12} - 1 = 4095$) cross-multiplied by
3 constraint dimensions, matching the theoretical prediction.

\subsection{Bitmask Operations in Hardware}

On FPGA, bitmask operations map to combinational logic with zero-cycle latency:
64-bit AND is a 2-LUT tree of depth 1 requiring 64 LUTs in Xilinx Artix-7
mapping. The FLUX hardware implementation dedicates 8 parallel 64-bit comparator
paths, each requiring 32 LUT6 primitives (2-input 64-bit equality), for a total
comparator cost of 256 LUTs---within the measured 1,717 LUT budget.

%% ============================================================
%% 4. FLUX ISA
%% ============================================================
\section{FLUX ISA: Constraint Enforcement Bytecode}
\label{sec:flux}

\subsection{Design Principles}

FLUX ISA is founded on four principles:

\begin{enumerate}
  \item \textbf{Constraint-first}: every tier includes the \texttt{Assert} opcode.
    A FLUX execution that omits constraint checking is architecturally non-compliant.
  \item \textbf{Tiered fidelity}: higher tiers are strict opcode supersets; the same
    \texttt{.flux} file is valid (though not fully executable) on any tier.
  \item \textbf{Stack-based execution}: no registers, no implicit side effects. The
    stack is the sole computational substrate, enabling formal stack-effect analysis.
  \item \textbf{Wire compatibility}: the \texttt{0x464C} (``FL'') binary format is
    shared across all tiers and hardware implementations.
\end{enumerate}

\subsection{Tier Architecture}

\begin{table}[h]
\centering
\caption{FLUX ISA implementation tiers.}
\label{tab:tiers}
\renewcommand{\arraystretch}{1.2}
\begin{tabular}{lllll}
\toprule
\textbf{Tier} & \textbf{Target} & \textbf{Opcodes} & \textbf{Runtime} & \textbf{Memory} \\
\midrule
\texttt{mini} & ARM Cortex-M0+/M4 & 21 (subset) & \texttt{\#![no\_std]} & 256 B stack \\
\texttt{std}  & Embedded Linux   & 35           & CLI toolchain     & 4 KB default \\
\texttt{edge} & Edge servers     & 35           & tokio async       & Dynamic + 256-slot \\
\texttt{thor} & GPU / fleet      & 43           & CUDA + tokio      & 65K stack + GPU \\
\bottomrule
\end{tabular}
\end{table}

\subsection{Wire Format}

The FLUX binary format uses a 4-byte header followed by a fixed-size instruction array:
\begin{verbatim}
  ┌──────────┬──────────────┬─────────────────────────┐
  │ Magic    │ Count (u16)  │ Instruction Array        │
  │ 0x464C   │ little-endian│ N × 24 bytes             │
  └──────────┴──────────────┴─────────────────────────┘

  Each 24-byte instruction:
  ┌────────┬──────┬──────────┬──────────┬──────────┬───────┐
  │ Opcode │ Pad  │ Operand0 │ Operand1 │ Reserved │ Flags │
  │ 1 byte │ 1 B  │ f64 8 B  │ f64 8 B  │ 4 bytes  │ u16   │
  └────────┴──────┴──────────┴──────────┴──────────┴───────┘
\end{verbatim}

Fixed 24-byte instruction size enables direct indexing without deserialization
overhead ($O(1)$ instruction fetch by index). We note an alignment hazard on
ARM Cortex-M: the 4-byte header places the instruction array at offset 4 from
buffer start, potentially misaligning \texttt{f64} fields requiring 8-byte
alignment. The \texttt{mini} tier provides an alignment-safe byte-copy decoder
for ARM targets.

\subsection{Opcode Taxonomy}

The 43 opcodes are organized in six functional groups:

\begin{table}[h]
\centering
\caption{FLUX ISA opcode taxonomy (selected opcodes).}
\label{tab:opcodes}
\renewcommand{\arraystretch}{1.15}
\begin{tabular}{lllll}
\toprule
\textbf{Group} & \textbf{Range} & \textbf{Key Opcodes} & \textbf{Stack Effect} \\
\midrule
Stack/Value  & 0x00--0x06 & \texttt{Push}, \texttt{Pop}, \texttt{Dup}, \texttt{Swap} & $\pm1$ \\
Arithmetic   & 0x10--0x15 & \texttt{Add}, \texttt{Sub}, \texttt{Mul}, \texttt{Div} & $-1$ \\
Logic/Compare& 0x20--0x35 & \texttt{And}, \texttt{Or}, \texttt{Eq}, \texttt{Lt} & $-1$ \\
Control Flow & 0x40--0x45 & \texttt{Jmp}, \texttt{Jz}, \texttt{Call}, \texttt{Halt} & varies \\
\textbf{CSP} & \textbf{0x50--0x54} & \textbf{\texttt{Assert}, \texttt{Constrain}, \texttt{Propagate}} & varies \\
Thor (GPU)   & 0x80--0x87 & \texttt{BatchSolve}, \texttt{GpuCompile}, \texttt{Pathfind} & varies \\
\bottomrule
\end{tabular}
\end{table}

\noindent The CSP group is the architectural core:

\begin{itemize}
  \item \texttt{Assert} (0x50): pops top of stack; halts with \texttt{ConstraintViolation}
    if zero. An \emph{irrecoverable gate}---execution cannot proceed past an \texttt{Assert}
    with the tested value on the stack.
  \item \texttt{Constrain} (0x51): pops value, lower bound, upper bound; pushes
    the domain-restricted value, establishing BitmaskDomain bounds.
  \item \texttt{Propagate} (0x52): executes arc-consistency reduction on the current
    domain state, implementing the REVISE algorithm in bytecode.
  \item \texttt{Solve} (0x53): on Thor, dispatches backtracking search to GPU.
  \item \texttt{Verify} (0x54): non-destructive domain membership check, leaving
    the value on the stack for downstream use.
\end{itemize}

\subsection{Constraint Semantics: Assert vs.\ Constrain}

The semantic distinction between \texttt{Assert} and \texttt{Constrain} enables
two constraint-programming styles:

\begin{enumerate}
  \item \textbf{Hard gates} (\texttt{Assert}): execution halts immediately on violation.
    Used for safety-critical invariants where operation beyond violation is inadmissible.
  \item \textbf{Soft domains} (\texttt{Constrain} + \texttt{Verify}): sets domain bounds
    without halting; downstream logic can query violation status. Used for constraint
    propagation within a search loop.
\end{enumerate}

\subsection{Static Validation}

Before execution, the FLUX bytecode validator performs two static checks:

\begin{enumerate}
  \item \textbf{Jump target validation}: all branch operands must reference valid instruction
    indices; out-of-bounds jumps are rejected at load time.
  \item \textbf{Stack effect analysis}: the validator tracks cumulative stack depth through
    the instruction stream, rejecting programs with statically detectable stack underflows.
\end{enumerate}

\begin{theorem}[Static Stack Soundness]
\label{thm:stack}
For any FLUX bytecode program $P$ that passes validation, no straight-line execution
path can produce a stack underflow.
\end{theorem}

\begin{proof}
The validator computes $\Delta_i = \text{outputs}_i - \text{inputs}_i$ for each
instruction $i$ and maintains a running cumulative sum $S_k = \sum_{i=0}^{k} \Delta_i$.
Validation rejects $P$ if $S_k < 0$ for any $k$. Therefore if validation passes,
$S_k \geq 0$ for all $k$---stack depth is non-negative throughout any straight-line
path. \qed
\end{proof}

\subsection{Compilation Pipeline: CSP to FLUX Bytecode}

The constraint compilation process transforms a declarative CSP specification (JSON
or GUARD) to FLUX bytecode through three phases:

\begin{description}
  \item[Phase 1---Domain Encoding] Each variable domain $D_i = [l_i, u_i]$ maps to a
    \texttt{Constrain} instruction sequence.
  \item[Phase 2---Constraint Flattening] Each constraint $c_j(x_{j_1}, x_{j_2})$ flattens
    to: push both variables, apply comparison, \texttt{Assert}. N-ary constraints decompose
    to binary combinations via \texttt{And}/\texttt{Or}.
  \item[Phase 3---Topology Ordering] Constraints are ordered to minimize stack depth;
    shared variables use \texttt{Dup}; live-variable analysis inserts \texttt{Pop} to
    release dead values.
\end{description}

\noindent For an 8-queens CSP ($n=8$, $\binom{8}{2}=28$ pairs, 2 constraints per pair),
the compiler produces 56 \texttt{Assert} blocks plus 8 \texttt{Constrain} instructions:

\begin{lstlisting}[style=rustcode,caption={8-Queens column-uniqueness constraint in FLUX bytecode.},label={lst:queens}]
; Constraint: q_i != q_j (column uniqueness)
PUSH qi_value   ; push column of queen i
PUSH qj_value   ; push column of queen j
Eq              ; check qi == qj
Not             ; invert: want qi != qj
Assert          ; halt if violated

; Constraint: |qi - qj| != |i - j| (diagonal)
PUSH qi_value
PUSH qj_value
Sub             ; compute qi - qj
PUSH offset     ; push |i - j|
Eq
Not
Assert          ; halt if diagonal violated
\end{lstlisting}

%% ============================================================
%% 5. GUARD DSL
%% ============================================================
\section{GUARD DSL: Declarative Constraint Specification}
\label{sec:guard}

\subsection{Design Philosophy}

GUARD (Generic Unified Assurance Requirement Descriptor) is a constraint
specification language designed for \emph{safety engineers}, not compiler hackers.
A GUARD module is simultaneously: a requirements document with formal semantics,
a FLUX compiler input, and a proof obligation generator. Six design principles
govern the language:

\begin{description}
  \item[P1---Readability] A safety engineer reads GUARD without training.
  \item[P2---Unit semantics] Physical units (\texttt{100 kt}, \texttt{2.5 g}) are
    semantic, not cosmetic; dimensional inconsistency is a type error.
  \item[P3---Explicit temporality] Temporal operators (\texttt{always P},
    \texttt{for 3s P}, \texttt{rate\_of x}) are first-class, not library functions.
  \item[P4---Single source of truth] The \texttt{.guard} file is the requirement;
    no separate Word document, no separate Simulink model.
  \item[P5---Proof independence] The compiler can be buggy; the proof certificate
    cannot lie, since the verifier ($\approx$500 lines of Rust) reruns all proof
    obligations independently.
  \item[P6---Graceful degradation] Properties are stratified as Proved (SMT
    solver), Bounded (bounded model checking), or Runtime (dynamic enforcement only).
\end{description}

\subsection{Language Structure}

A GUARD module has seven sections:

\begin{lstlisting}[style=guardcode,caption={GUARD module structure for a throttle controller.},label={lst:guard}]
module ThrottleLimit version "1.0.0"
  target FluxISA-43

dimension Percent is real from 0 % to 100 %

state throttle_command has Percent
state engine_enabled   has bool

invariant ThrottleMustNotExceedMax
  critical
  ensure throttle_command <= 100 %
  on_violation halt;

invariant ThrottleMustNotReverse
  critical
  ensure throttle_command >= 0 %
  on_violation halt;

derive PositiveThrottleImpliesEngineEnabled
  from throttle_command > 0 %
  conclude engine_enabled = true

proof ThrottleSafetyProof
  tactic interval_arithmetic
  tactic bit_blast
\end{lstlisting}

\subsection{Temporal Semantics}

GUARD uses discrete-time LTL with bounded-duration operators over a trace
$\sigma_0, \sigma_1, \sigma_2, \ldots$ of system states:

\begin{align*}
  \texttt{always } P &\equiv \forall k \geq 0.\; P(\sigma_k) \\
  \texttt{eventually } P &\equiv \exists k \geq 0.\; P(\sigma_k) \\
  \texttt{for } T\; P &\equiv \forall k \in [t, t{+}n].\; P(\sigma_k),\; T = n \cdot \Delta t \\
  \texttt{rate\_of } x &\equiv (\sigma_t(x) - \sigma_{t-1}(x)) / \Delta t
\end{align*}

For each invariant $I$, the compiler generates two verification conditions:
\begin{enumerate}
  \item \textbf{Initiation}: $\text{InitialState} \Rightarrow I(\sigma_0)$
  \item \textbf{Preservation}: $I(\sigma_t) \wedge \text{Trans}(\sigma_t, \sigma_{t+1}) \Rightarrow I(\sigma_{t+1})$
\end{enumerate}
Proving both establishes $I$ as an inductive invariant holding in all reachable states.

\subsection{Proof Certificates}

Every compiled GUARD module co-generates a \texttt{.guardcert} JSON artifact:
\begin{itemize}
  \item SHA-256 hashes of source and bytecode (tamper detection).
  \item Per-invariant verification conditions in SMT-LIB format.
  \item Solver results (\texttt{sat}/\texttt{unsat}/\texttt{unknown}) with
    counterexamples for failed proofs.
  \item Merkle root over all obligation trace digests.
\end{itemize}

The verifier's Trusted Computing Base (TCB) is minimal:
(1)~SHA-256, (2)~an SMT solver (Z3/cvc5/Bitwuzla),
(3)~the FLUX VM ($\approx$300 lines), (4)~the GUARD source parser ($\approx$200 lines).
The compiler itself is \emph{not} in the TCB; a buggy compiler cannot forge a
valid certificate because the verifier re-runs all proof obligations independently.

The Merkle construction provides a tamper-evident seal:
\[
  \ell_i = \mathsf{SHA256}(\mathsf{id}_i \| \mathsf{trace\_digest}_i), \quad
  \mathsf{root} = \mathsf{reduce}(\{\ell_i\},\; \lambda\,a,b.\, \mathsf{SHA256}(a\|b))
\]
If any obligation is recomputed with different solver parameters, its trace digest
changes, invalidating the root. Auditors detect regressions with a single hash comparison.

%% ============================================================
%% 6. DO-254 DAL A HARDWARE
%% ============================================================
\section{DO-254 DAL A Certification Path}
\label{sec:hardware}

\subsection{Hardware Architecture}

The FLUX constraint checker is implemented as a SystemVerilog module targeting
Xilinx Artix-7 (xc7a100t) synthesized at 100 MHz. Figure~\ref{fig:hardware} shows
the top-level architecture. Safety-critical design techniques follow
DO-254~\cite{do254} \S6.3.2 and the ARP4754A joint safety argument:

\begin{figure}[h]
\centering
\begin{tikzpicture}[node distance=1.2cm, auto,
  block/.style={rectangle, draw, fill=blue!10, text width=3.2cm,
                text centered, rounded corners, minimum height=0.9cm},
  fault/.style={rectangle, draw, fill=red!20, text width=3.2cm,
                text centered, rounded corners, minimum height=0.9cm}]

  \node [block] (dec)  {Opcode Decoder\\(1-hot, 43 opcodes)};
  \node [block, below of=dec] (stack) {Parity Stack\\(64-bit, depth 16)};
  \node [block, below of=stack] (comp) {Parallel Comparators\\(8×64-bit)};
  \node [fault, right of=dec, node distance=4.5cm] (tmr) {TMR Registers\\(PC, SP, FSM)};
  \node [fault, below of=tmr] (latch) {Fault Latch\\(set-once, no clear)};
  \node [fault, below of=latch] (ssc) {Safe-State Controller\\(3-state FSM)};

  \draw[->] (dec) -- node[right]{opcode\_hot} (stack);
  \draw[->] (stack) -- (comp);
  \draw[->] (tmr) -- node[right]{mismatch} (latch);
  \draw[->] (dec) -- node[above]{illegal\_op} (latch);
  \draw[->] (stack) -- node[above]{bounds/parity} (latch);
  \draw[->] (latch) -- node[right]{global\_fault} (ssc);
  \draw[->] (ssc) -- node[below]{vm\_nop\_force} (comp);
\end{tikzpicture}
\caption{FLUX constraint checker hardware architecture.}
\label{fig:hardware}
\end{figure}

\subsection{Triple Modular Redundancy (TMR)}

All safety-critical control state (program counter, stack pointer, FSM state)
is stored in explicit triple registers with majority voting:
\[
  q_i = (a_i \wedge b_i) \vee (b_i \wedge c_i) \vee (a_i \wedge c_i)
\]
Synthesis is protected by \texttt{keep = "true", dont\_touch = "true"} attributes
preventing optimizer merging of replicas. TMR is applied to control registers only
(partial TMR); DO-254 \S6.3.2.2 Note~3 permits this when residual risk of data-path
single-event upsets is shown to be acceptably low.

\subsection{Parity-Protected Stack}

Each 64-bit stack entry is stored with a parity bit:
\[
  \mathsf{parity}(d) = \bigoplus_{i=0}^{63} d_i
\]
Stack write appends the parity bit; stack read checks parity, asserting
\texttt{FAULT\_STACK\_PARITY} if corrupted. Bounds checking detects push-on-full
and pop-on-empty.

\subsection{One-Hot Opcode Decoder}

The opcode decoder maps 6-bit opcode to a 43-bit one-hot vector using a
combinational \texttt{unique case}. An opcode is \emph{illegal} if and only if
$\mathsf{POPCOUNT}(\mathsf{opcode\_hot}) \neq 1$:
\begin{lstlisting}[language=Verilog,style=rustcode,caption={One-hot decode with illegal-op detection.},label={lst:onehot}]
assign illegal_op_fault = $countones(opcode_onehot) != 6'd1;
\end{lstlisting}

\subsection{Safe-State Controller}

The safe-state controller implements a latching FSM:
$\mathsf{RUN} \to \mathsf{INTERLOCK} \to \mathsf{SAFE}$. Once in
\textsc{Safe}, the state is permanent until power cycle. Safe-state actions:
\begin{itemize}
  \item Force all comparator outputs to zero (safe-value assertion).
  \item Force the VM to execute \texttt{NOP} on every cycle.
  \item Assert the interlock output pin.
  \item Toggle the status pin at 1 Hz heartbeat.
\end{itemize}

This design reflects a key insight from multi-model review: \textbf{HALT is not a
safe state}. A constraint violation must transition the hardware to a \emph{defined
hazard-free state}, not merely cease execution~\cite{deep-research-2026}.

\subsection{FPGA Resource Utilization}

\begin{table}[h]
\centering
\caption{FPGA synthesis results: Xilinx Artix-7 xc7a100t, 100 MHz.}
\label{tab:fpga}
\renewcommand{\arraystretch}{1.2}
\begin{tabular}{lrrr}
\toprule
\textbf{Resource} & \textbf{Used} & \textbf{Available} & \textbf{Utilization} \\
\midrule
LUT6              & 1,717  & 101,400 & 1.69\% \\
Flip-Flops        & 1,807  & 202,800 & 0.89\% \\
BRAM              & 0      & 135     & 0.00\% \\
DSP48             & 0      & 240     & 0.00\% \\
\midrule
\multicolumn{4}{l}{\textit{Power breakdown}} \\
\midrule
Static power      & \multicolumn{3}{l}{85 mW} \\
Dynamic power     & \multicolumn{3}{l}{35 mW} \\
\textbf{Total power}       & \multicolumn{3}{l}{\textbf{120 mW}} \\
\midrule
Violation detection latency & \multicolumn{3}{l}{$\leq$2 clock cycles (20 ns at 100 MHz)} \\
\bottomrule
\end{tabular}
\end{table}

The 1,717-LUT footprint leaves 98.3\% of the Artix-7 100T available for
application logic, enabling the FLUX checker to operate as a dedicated safety
co-processor on resource-constrained flight controllers.

\subsection{Formal Verification of Hardware Safety Properties}

Seven SymbiYosys/Z3 properties are verified over the RTL (12-step bounded model
checking with proof extension of 4):

\begin{table}[h]
\centering
\caption{Formally verified safety properties.}
\label{tab:formal}
\renewcommand{\arraystretch}{1.2}
\begin{tabular}{lll}
\toprule
\textbf{Property} & \textbf{Type} & \textbf{Formal Statement} \\
\midrule
Fault sticky & Safety & \texttt{fault |=> always fault} \\
Interlock latency & Safety & \texttt{\$rose(fault) |-> interlock\_out} \\
Safe NOP force & Safety & \texttt{fault |-> vm\_nop\_force} \\
TMR detection & Safety & \texttt{tmr.mismatch |-> fault} \\
Stack bounds & Safety & \texttt{stack.bounds\_fault |-> \#\#1 fault} \\
Illegal opcode & Safety & \texttt{dec.illegal\_op |-> \#\#1 fault} \\
Safe pin active & Safety & \texttt{fault |-> \#\#2 safe\_pin != z} \\
\bottomrule
\end{tabular}
\end{table}

All seven properties are proved; no counterexamples found within the 12-step depth.
Coverage properties additionally confirm that normal execution and fault-entry
states are reachable.

%% ============================================================
%% 7. COQ FORMALIZATION
%% ============================================================
\section{Coq Formalization: P2 Invariant Preservation}
\label{sec:coq}

\subsection{Setup}

We formalize the arc-consistency engine in Coq using propositional domains.
A domain is a predicate $D_i : \mathbb{N} \to \mathsf{Prop}$; a constraint is a
binary relation $c : \mathbb{N} \to \mathbb{N} \to \mathsf{Prop}$.

\begin{lstlisting}[style=coqcode,caption={Core domain and constraint types.},label={lst:coq-types}]
Definition Domain := Val -> Prop.
Definition Constraint := Val -> Val -> Prop.

Definition INV (D : Var -> Domain) (arcs : list Arc) : Prop :=
  exists alpha : Assignment,
    within_domains alpha D /\
    satisfies_all alpha arcs.

Definition revise (Di Dj : Domain) (c : Constraint) : Domain :=
  fun x => Di x /\ exists y, Dj y /\ c x y.
\end{lstlisting}

\noindent \texttt{INV} asserts that a satisfying assignment exists within the
current domains. \texttt{revise} implements the REVISE operation: retaining only
values in $D_i$ that have at least one support in $D_j$ under constraint $c$.

\subsection{Three Proved Theorems}

\begin{theorem}[REVISE Preserves INV]
\label{thm:revise}
Let $D$ be a domain function and $\mathit{arcs}$ a set of constraints such that
$\mathsf{INV}(D, \mathit{arcs})$ holds. Let $a \in \mathit{arcs}$ and let
$D'_i = \mathsf{revise}(D_i, D_j, c_a)$. Then
$\mathsf{INV}(\mathsf{update}(D, i, D'_i), \mathit{arcs})$.
\end{theorem}

\begin{proof}
We use the witness $\alpha$ from $\mathsf{INV}(D, \mathit{arcs})$. We must show
$\alpha$ still satisfies all domains in the updated system. For variables $k \neq i$,
the domain is unchanged so $\alpha(k) \in D_k$ by hypothesis. For variable $i$,
we need $\alpha(i) \in D'_i = \mathsf{revise}(D_i, D_j, c_a)$. Since $\alpha$
satisfies arc $a$, $c_a(\alpha(i), \alpha(j))$ holds, providing the required support
$y = \alpha(j) \in D_j$. Therefore $\alpha(i) \in D'_i$. Since $\alpha$ satisfies
all arcs in $\mathit{arcs}$ by hypothesis, the updated system is still INV. \qed
\end{proof}

\begin{theorem}[Arc-Consistency Propagation Preserves INV]
\label{thm:union}
For any sequence of arcs $\mathit{worklist} \subseteq \mathit{arcs}$ and any $D$
with $\mathsf{INV}(D, \mathit{arcs})$, applying $\mathsf{revise}$ to each arc in
$\mathit{worklist}$ preserves $\mathsf{INV}$.
\end{theorem}

\begin{proof}
By induction on the worklist. Base case: empty worklist, $D$ unchanged, trivially
$\mathsf{INV}$. Inductive step: by Theorem~\ref{thm:revise}, processing one arc
preserves $\mathsf{INV}$; the induction hypothesis applies to the tail. \qed
\end{proof}

\begin{theorem}[HALT Means Not-INV]
\label{thm:halt}
If any variable domain $D_i$ is empty (the system halts), then
$\neg\, \mathsf{INV}(D, \mathit{arcs})$.
\end{theorem}

\begin{proof}
Suppose $\mathsf{INV}(D, \mathit{arcs})$ holds; then there exists a satisfying
assignment $\alpha$ with $\alpha(i) \in D_i$ for all $i$. But $D_i = \emptyset$
means no such $\alpha(i)$ exists---contradiction. Therefore $\neg\, \mathsf{INV}$. \qed
\end{proof}

\subsection{Summary of Guarantees}

Together the three theorems establish:
\begin{itemize}
  \item \textbf{Soundness}: the propagation engine never eliminates a globally satisfying
    solution (Theorems~\ref{thm:revise}, \ref{thm:union}).
  \item \textbf{Halt completeness}: when the FLUX machine detects an empty domain and
    asserts \textsc{Halt}, the CSP is provably unsatisfiable under the current domains
    (Theorem~\ref{thm:halt}).
\end{itemize}

The Coq file (\texttt{flux\_p2.v}, $\approx$325 lines) additionally proves bitmask
correctness lemmas: \texttt{bitmask\_and\_is\_inter} and \texttt{bitmask\_andnot\_is\_diff}
establish that bitwise AND and ANDNOT are correct implementations of domain intersection
and domain difference respectively, connecting the abstract Coq model to the hardware
implementation.

\subsection{Open Gap: Compilation Correctness}

We cannot currently prove \emph{compilation correctness}---that the FLUX bytecode
emitted from a CSP specification is a faithful representation of the original constraints.
This requires a formal correspondence between the CSP specification language and FLUX
bytecode semantics. We term this the \emph{Lean 4 gap}: it is planned as the
highest-priority future formalization effort, following the approach of
CompCert~\cite{leroy2009} for compiled constraint programs.

%% ============================================================
%% 8. EVALUATION
%% ============================================================
\section{Evaluation}
\label{sec:eval}

\subsection{Benchmark Suite}

We evaluate on three benchmark problems representing practical constraint domains:

\begin{enumerate}
  \item \textbf{N-Queens}: classical combinatorial CSP; $n$ queens on an $n \times n$
    board with all-different row/column/diagonal constraints.
  \item \textbf{Graph Coloring}: chromatic number problem on random $G(n,p)$ graphs;
    $k$-coloring with adjacent-vertex inequality constraints.
  \item \textbf{Job-Shop Scheduling}: $m$ machines, $n$ jobs; no-overlap constraints
    on machine assignment windows.
\end{enumerate}

\subsection{BitmaskDomain vs.\ General Solver}

\begin{table}[h]
\centering
\caption{N-Queens: BitmaskDomain vs.\ list-based solver (x86-64, Ryzen 9 5900X).}
\label{tab:nqueens}
\renewcommand{\arraystretch}{1.2}
\begin{tabular}{rrrrl}
\toprule
\textbf{N} & \textbf{Bitmask ($\mu$s)} & \textbf{General ($\mu$s)} & \textbf{Solutions} & \textbf{Speedup} \\
\midrule
4  & $<$1    & $<$1    & 2         & ---      \\
6  & $<$1    & $<$1    & 4         & ---      \\
8  & $<$1    & 84      & 92        & $>84\times$ \\
10 & 12      & 1,847   & 724       & $154\times$ \\
12 & 180     & 48,210  & 14,200    & $268\times$ \\
14 & 3,420   & ---     & 365,596   & ---      \\
\bottomrule
\end{tabular}
\end{table}

\begin{table}[h]
\centering
\caption{Graph coloring: BitmaskDomain speedup on random $G(50, 0.3)$ instances.}
\label{tab:coloring}
\renewcommand{\arraystretch}{1.2}
\begin{tabular}{lrrr}
\toprule
\textbf{Colors} & \textbf{Bitmask ($\mu$s)} & \textbf{General ($\mu$s)} & \textbf{Speedup} \\
\midrule
$k=3$ (sat) & 210  & 14,800  & $70\times$ \\
$k=4$ (sat) & 18   & 830     & $46\times$ \\
$k=5$ (sat) & 3    & 95      & $32\times$ \\
Unsat detection & 4,100 & ---  & ---    \\
\bottomrule
\end{tabular}
\end{table}

\begin{table}[h]
\centering
\caption{Job-shop scheduling: $6\times6$ instances (6 jobs, 6 machines).}
\label{tab:jobshop}
\renewcommand{\arraystretch}{1.2}
\begin{tabular}{lrr}
\toprule
\textbf{Instance} & \textbf{Bitmask ($\mu$s)} & \textbf{General ($\mu$s)} \\
\midrule
FT06 (optimal 55) & 340  & 41,200 \\
LA01 (optimal 666)& 2,800 & 390,000 \\
Random tight & 190  & 17,600 \\
\bottomrule
\end{tabular}
\end{table}

\noindent The speedup grows with domain cardinality, consistent with
Theorem~\ref{thm:speedup}. The 12-Queens instance shows 268$\times$ speedup;
larger domains approach the $12{,}324\times$ theoretical ceiling.

\subsection{FLUX Bytecode Execution: Sonar Physics}

As a domain-specific case study, we compile Mackenzie~1981 sound speed bounds
and Fran\c{c}ois-Garrison~1982 absorption constraints~\cite{francois1982}
to FLUX bytecode for the \texttt{mini} tier (ARM Cortex-M4F):

\begin{table}[h]
\centering
\caption{FLUX constraint check timing on ARM Cortex-M4F (STM32F407, 168 MHz).}
\label{tab:arm}
\renewcommand{\arraystretch}{1.2}
\begin{tabular}{lrr}
\toprule
\textbf{Operation} & \textbf{Cycles (est.)} & \textbf{Time ($\mu$s)} \\
\midrule
\texttt{Assert} (integer compare)  & 10--15      & $<0.1$ \\
\texttt{Add} (f64 via softfp)      & 30--50      & 0.18--0.30 \\
Sonar physics check (5 constraints) & 200--400  & 1.2--2.4 \\
Full Mackenzie validation (9-term) & 2,000--4,000 & 12--24 \\
\bottomrule
\end{tabular}
\end{table}

Sonar sampling intervals are typically 1--100 ms; the 24 $\mu$s worst-case
constraint check occupies $<$2.4\% of the sampling budget at maximum check
frequency, confirming real-time feasibility.

\subsection{GPU Batch Solving (Thor Tier)}

The Thor-tier batch solver dispatches CSP instances to GPU when batch size
exceeds 256 instances:

\begin{table}[h]
\centering
\caption{Thor-tier batch CSP solving: CPU (rayon) vs.\ GPU (CUDA, projected).}
\label{tab:gpu}
\renewcommand{\arraystretch}{1.2}
\begin{tabular}{rrrr}
\toprule
\textbf{Batch Size} & \textbf{CPU ($\mu$s)} & \textbf{GPU ($\mu$s)} & \textbf{GPU Speedup} \\
\midrule
10      & $\sim$50   & $\sim$200* & 0.25$\times$ \\
100     & $\sim$500  & $\sim$250  & 2$\times$ \\
1,000   & $\sim$5,000 & $\sim$400 & 12.5$\times$ \\
10,000  & $\sim$50,000 & $\sim$1,500 & 33$\times$ \\
\bottomrule
\multicolumn{4}{l}{\footnotesize *GPU dispatch overhead dominates for small batches.} \\
\multicolumn{4}{l}{\footnotesize GPU numbers are projected; CUDA kernel compilation in progress.}
\end{tabular}
\end{table}

\subsection{Comparison with Prior Systems}

\begin{table*}[t]
\centering
\caption{Comparison of FLUX with related constraint enforcement systems.}
\label{tab:comparison}
\renewcommand{\arraystretch}{1.2}
\begin{tabular}{lccccccc}
\toprule
\textbf{Property} & \textbf{FLUX ISA} & \textbf{Eiffel/Spec\#} & \textbf{Bitcoin Script} & \textbf{eBPF} & \textbf{ROS~2} & \textbf{TLA+} & \textbf{MISRA~C} \\
\midrule
Constraints as opcodes    & \checkmark & $\times$ (lang-level) & \checkmark & $\times$ (verifier) & $\times$ & $\times$ & $\times$ \\
Portable bytecode         & \checkmark & $\times$ & \checkmark & \checkmark & $\times$ & $\times$ & $\times$ \\
Runs on MCU               & \checkmark & $\times$ & $\times$ & $\times$ & $\times$ & $\times$ & \checkmark \\
GPU acceleration          & \checkmark & $\times$ & $\times$ & $\times$ & $\times$ & $\times$ & $\times$ \\
Domain physics constraints& \checkmark & \checkmark & $\times$ & $\times$ & $\times$ & \checkmark & $\times$ \\
DO-254 certified hardware & \checkmark & $\times$ & $\times$ & $\times$ & $\times$ & $\times$ & $\times$ \\
Proof certificates        & \checkmark & Partial & $\times$ & \checkmark & $\times$ & \checkmark & $\times$ \\
Formal compilation proof  & Partial & Partial & $\times$ & \checkmark & $\times$ & \checkmark & $\times$ \\
Fleet coordination        & \checkmark & $\times$ & $\times$ & $\times$ & \checkmark & $\times$ & $\times$ \\
\bottomrule
\end{tabular}
\end{table*}

%% ============================================================
%% 9. THREATS TO VALIDITY
%% ============================================================
\section{Threats to Validity}
\label{sec:threats}

\textbf{Benchmark representativeness.} N-Queens is a classical benchmark but is not
representative of industrial constraint problems. We include graph coloring and
job-shop scheduling to broaden coverage; real avionics and automotive CSPs are left
for future industrial evaluation.

\textbf{GPU numbers are projected.} The Thor-tier GPU batch numbers in
Table~\ref{tab:gpu} are CPU measurements extrapolated to GPU; the CUDA kernel
compilation path is under development. Hardware validation will replace projections
in the full version.

\textbf{ARM timing is estimated.} Cortex-M4F constraint check timing
(Table~\ref{tab:arm}) is derived from instruction-count analysis, not DWT cycle
counter measurements on real hardware. We plan to report measured results.

\textbf{Compilation correctness gap.} The Coq formalization proves properties of the
VM semantics and propagation engine, but does not yet prove that the GUARD/JSON
compiler produces bytecode that faithfully represents the input specification. This
gap is clearly identified; closing it is the highest-priority future work.

\textbf{SEU vulnerability.} A single-event upset (SEU) in FPGA configuration memory
can rewrite FLUX bytecode, bypassing all safety properties. Production deployments
require configuration memory scrubbing or radiation-hardened FPGAs. This is a known
limitation of commercial SRAM-based FPGAs in high-reliability environments.

%% ============================================================
%% 10. CONCLUSION
%% ============================================================
\section{Conclusion}
\label{sec:conclusion}

We have presented constraint-native computing: an architecture in which constraint
satisfaction is enforced as a structural property of bytecode execution and certified
silicon, not as an optional runtime check.

The central insight is both simple and practical: \textbf{if wrong answers are
possible, wrong should halt execution}. FLUX ISA makes this architectural, not
aspirational. The \texttt{Assert} opcode occupies the same position in the
instruction hierarchy as \texttt{Add}---it cannot be bypassed, disabled in
production, or forgotten by a software engineer.

The contributions compose into a coherent certification stack:

\begin{enumerate}
  \item \textbf{BitmaskDomain} makes constraint solving fast enough for real-time
    embedded deployment: single-instruction domain operations with speedup scaling
    linearly with domain cardinality.

  \item \textbf{FLUX ISA} provides a portable, tiered bytecode substrate where
    constraint checking spans four orders of magnitude in hardware capability, from
    a 256-byte ARM microcontroller to a GPU fleet, sharing the same wire format
    and the same constraint semantics.

  \item \textbf{GUARD DSL} bridges the semantic gap at the specification level:
    safety engineers write readable requirements; the compiler generates both
    certified bytecode and machine-checkable proof certificates with a minimal
    trusted computing base.

  \item \textbf{DO-254 DAL~A hardware} demonstrates that programmable constraint
    enforcement fits within avionics-grade FPGA resource budgets (1,717 LUTs,
    120 mW) with deterministic violation detection latency ($\leq$20 ns).

  \item \textbf{Coq mechanization} provides a machine-verified foundation for the
    arc-consistency engine: propagation is sound (never eliminates solutions); halt
    is complete (empty domain certifies unsatisfiability).
\end{enumerate}

\noindent The frontier left open is compilation correctness: proving that GUARD
source files compile to bytecode programs that faithfully implement their specifications.
This is the analogue of CompCert's simulation theorem for the constraint compilation
domain, and we believe it is tractable using the same $k$-induction and refinement
techniques applied in the GUARD proof certificates.

Constraint-native computing is not a new paradigm invented from whole cloth.
The lineage from Eiffel's Design by Contract through Bitcoin Script's
\texttt{OP\_VERIFY} through eBPF's load-time verifier demonstrates that building
correctness into the execution layer is a recurring, productive architectural idea.
FLUX ISA's contribution is applying this idea to a new problem domain---physical
constraints in deployed autonomous systems---with a new implementation approach:
tiered bytecode running from microcontrollers to GPU clusters, synthesized to
certified hardware, and mechanically verified in Coq.

The constraint \emph{is} the hardware. The hardware \emph{is} the constraint.

%% ============================================================
%% ACKNOWLEDGMENTS
%% ============================================================
\begin{acks}
The authors thank the Cocapn Fleet for review, the multi-model synthesis sessions
with Claude Opus~4, Kimi, and Nemotron-3-Super that sharpened the hardware
threat model, and Seed-2.0-pro for providing the first real Vivado synthesis numbers.
\end{acks}

%% ============================================================
%% REFERENCES
%% ============================================================
\bibliographystyle{ACM-Reference-Format}
\begin{thebibliography}{99}

\bibitem{barnett2011}
M.~Barnett, M.~F{\"a}hndrich, K.~R.~M.~Leino, P.~M{\"u}ller, W.~Schulte, and H.~Venter.
\newblock Specification and verification: The {Spec\#} experience.
\newblock {\em Communications of the ACM}, 54(10):81--91, 2011.

\bibitem{demoura2021}
L.~de~Moura and S.~Ullrich.
\newblock The {Lean~4} theorem prover and programming language.
\newblock In {\em CADE-28}, 2021.

\bibitem{do254}
{RTCA/DO-254}.
\newblock Design assurance guidance for airborne electronic hardware.
\newblock Technical report, RTCA, 2000.

\bibitem{francois1982}
R.~E.~Fran\c{c}ois and G.~R.~Garrison.
\newblock Sound absorption based on ocean measurements: Part {II}: Boric acid contribution and equation for total absorption.
\newblock {\em Journal of the Acoustical Society of America}, 72(6):1879--1890, 1982.

\bibitem{gregg2019}
B.~Gregg.
\newblock {\em {BPF} Performance Tools}.
\newblock Addison-Wesley, 2019.

\bibitem{haas2017}
A.~Haas et~al.
\newblock Bringing the web up to speed with {WebAssembly}.
\newblock In {\em PLDI~'17}, pages 185--200, 2017.

\bibitem{halbwachs1991}
N.~Halbwachs, P.~Caspi, P.~Raymond, and D.~Pilaud.
\newblock The synchronous data flow programming language {LUSTRE}.
\newblock {\em Proceedings of the IEEE}, 79(9):1305--1320, 1991.

\bibitem{jackson2012}
D.~Jackson.
\newblock {\em Software Abstractions: Logic, Language, and Analysis}.
\newblock MIT Press, 2012.

\bibitem{lamport2002}
L.~Lamport.
\newblock {\em Specifying Systems: The {TLA+} Language and Tools for Hardware and Software Engineers}.
\newblock Addison-Wesley, 2002.

\bibitem{leroy2009}
X.~Leroy.
\newblock Formal verification of a realistic compiler.
\newblock {\em Communications of the ACM}, 52(7):107--115, 2009.

\bibitem{mackenzie1981}
K.~V.~Mackenzie.
\newblock Nine-term equation for sound speed in the oceans.
\newblock {\em Journal of the Acoustical Society of America}, 70(3):807--812, 1981.

\bibitem{macenski2022}
S.~Macenski et~al.
\newblock The {Marathon~2}: A system using the {ROS~2} navigation stack.
\newblock {\em IEEE Robotics and Automation Magazine}, 2022.

\bibitem{meyer1992}
B.~Meyer.
\newblock Applying ``design by contract.''
\newblock {\em Computer}, 25(10):40--51, 1992.

\bibitem{misra2012}
{MISRA~C:2012}.
\newblock Guidelines for the use of the {C} language in critical systems.
\newblock Technical report, Motor Industry Software Reliability Association, 2012.

\bibitem{nakamoto2008}
S.~Nakamoto.
\newblock Bitcoin: A peer-to-peer electronic cash system.
\newblock Technical report, 2008.

\bibitem{perron2019}
L.~Perron and V.~Furnon.
\newblock {OR-Tools}.
\newblock \url{https://developers.google.com/optimization}, 2019.

\bibitem{schulte2009}
C.~Schulte and G.~Tack.
\newblock {\em Gecode: A Generic Constraint Development Environment}.
\newblock \url{https://www.gecode.org/}, 2009.

\bibitem{deep-research-2026}
Forgemaster~{\ding{42}}.
\newblock Multi-model strategic analysis: Constraint-native computing.
\newblock Technical report, Cocapn Fleet, 2026-05-03.

\end{thebibliography}

%% ============================================================
%% APPENDIX A: Complete Opcode Reference
%% ============================================================
\appendix
\section{Complete FLUX Opcode Reference}
\label{app:opcodes}

\begin{table}[h]
\centering
\caption{Base opcode set (35 opcodes, all tiers).}
\label{tab:opcodes-full}
\renewcommand{\arraystretch}{1.1}
\begin{tabular}{cllll}
\toprule
\textbf{Hex} & \textbf{Mnemonic} & \textbf{Group} & \textbf{Stack} & \textbf{Description} \\
\midrule
0x00 & \texttt{Nop}       & Stack    & $\pm0$ & No operation \\
0x01 & \texttt{Push}      & Stack    & $+1$   & Push immediate \\
0x02 & \texttt{Pop}       & Stack    & $-1$   & Discard top \\
0x03 & \texttt{Dup}       & Stack    & $+1$   & Duplicate top \\
0x04 & \texttt{Swap}      & Stack    & $\pm0$ & Exchange top two \\
0x05 & \texttt{Load}      & Stack    & $+1$   & Load from slot \\
0x06 & \texttt{Store}     & Stack    & $-2$   & Store to slot \\
\midrule
0x10 & \texttt{Add}       & Arith.   & $-1$   & $a + b$ \\
0x11 & \texttt{Sub}       & Arith.   & $-1$   & $a - b$ \\
0x12 & \texttt{Mul}       & Arith.   & $-1$   & $a \times b$ \\
0x13 & \texttt{Div}       & Arith.   & $-1$   & $a / b$ (halt on ÷0) \\
0x14 & \texttt{Mod}       & Arith.   & $-1$   & $a \bmod b$ \\
0x15 & \texttt{Neg}       & Arith.   & $\pm0$ & Negate \\
\midrule
0x20 & \texttt{And}       & Logic    & $-1$   & Logical AND \\
0x21 & \texttt{Or}        & Logic    & $-1$   & Logical OR \\
0x22 & \texttt{Not}       & Logic    & $\pm0$ & Logical NOT \\
0x30 & \texttt{Eq}        & Compare  & $-1$   & $a = b$? \\
0x31 & \texttt{Ne}        & Compare  & $-1$   & $a \neq b$? \\
0x32 & \texttt{Lt}        & Compare  & $-1$   & $a < b$? \\
0x33 & \texttt{Le}        & Compare  & $-1$   & $a \leq b$? \\
0x34 & \texttt{Gt}        & Compare  & $-1$   & $a > b$? \\
0x35 & \texttt{Ge}        & Compare  & $-1$   & $a \geq b$? \\
\midrule
0x40 & \texttt{Jmp}       & Control  & $+0$   & Unconditional jump \\
0x41 & \texttt{Jz}        & Control  & $-1$   & Jump if false \\
0x42 & \texttt{Jnz}       & Control  & $-1$   & Jump if true \\
0x43 & \texttt{Call}      & Control  & $+0$   & Call subroutine \\
0x44 & \texttt{Ret}       & Control  & $+0$   & Return \\
0x45 & \texttt{Halt}      & Control  & $+0$   & Stop execution \\
\midrule
\textbf{0x50} & \textbf{\texttt{Assert}}    & \textbf{CSP} & $-1$   & \textbf{Halt if false} \\
0x51 & \texttt{Constrain} & CSP      & $-3/+1$ & Set domain bounds \\
0x52 & \texttt{Propagate} & CSP      & $+1$   & Arc consistency \\
0x53 & \texttt{Solve}     & CSP      & $+1$   & CSP solve (GPU) \\
0x54 & \texttt{Verify}    & CSP      & $\pm0$ & Non-destructive check \\
\midrule
0x60 & \texttt{Print}     & I/O      & $+0$   & Debug output \\
0x61 & \texttt{Debug}     & I/O      & $+0$   & Trace output \\
\bottomrule
\end{tabular}
\end{table}

\begin{table}[h]
\centering
\caption{Thor extended opcodes (8 opcodes, thor tier only).}
\label{tab:thor-opcodes}
\renewcommand{\arraystretch}{1.1}
\begin{tabular}{cllll}
\toprule
\textbf{Hex} & \textbf{Mnemonic} & \textbf{Stack} & \textbf{Description} \\
\midrule
0x80 & \texttt{ParallelBranch} & $+1$    & Spawn $N$ tokio tasks \\
0x81 & \texttt{Reduce}         & $-1$    & Merge parallel results \\
0x82 & \texttt{GpuCompile}     & $+1$    & Compile to CUDA kernel \\
0x83 & \texttt{BatchSolve}     & $-1/+1$ & Batch CSP on GPU \\
0x84 & \texttt{SonarBatch}     & $-1/+1$ & Sonar physics GPU batch \\
0x85 & \texttt{TileCommit}     & $+0$    & Commit to PLATO graph \\
0x86 & \texttt{Pathfind}       & $+1$    & Traverse PLATO graph \\
0x87 & \texttt{ExtendedEnd}    & $+0$    & End extension sequence \\
\bottomrule
\end{tabular}
\end{table}

%% ============================================================
%% APPENDIX B: GUARD Flight Envelope Example
%% ============================================================
\section{GUARD Flight Envelope Example}
\label{app:flight}

\begin{lstlisting}[style=guardcode,caption={GUARD flight envelope constraint (excerpt).},label={lst:flight}]
module FlightEnvelope version "1.0.0"
  target FluxISA-43

dimension Knots   is real from 0 kt to 500 kt
dimension Feet    is real from -500 ft to 60000 ft
dimension G_Force is real from -3 g to 10 g

state airspeed  has Knots
state altitude  has Feet
state load_factor has G_Force

invariant VneNotExceeded        -- velocity never exceed
  critical
  ensure airspeed <= 340 kt
  on_violation halt;

invariant AltitudeSafeRange
  critical
  ensure altitude >= 0 ft and altitude <= 45000 ft
  on_violation halt;

invariant LoadFactorSafe
  critical
  when airspeed > 250 kt
  ensure load_factor <= 3.5 g and load_factor >= -1.5 g
  on_violation halt;

invariant StallSpeedRespected
  critical
  ensure not (airspeed < 120 kt and altitude < 1000 ft)
  on_violation halt;

invariant GForceRateLimit       -- temporal: G onset
  ensure always
    (rate_of load_factor <= 6 g / 1 s)
  on_violation halt;

proof FlightSafetyProof
  tactic k_induction 5
  tactic interval_arithmetic
\end{lstlisting}

\end{document}
